\documentclass[12pt]{article}
\usepackage{amsmath,amssymb,amsfonts}
\usepackage[margin=1in]{geometry}
\begin{document}
\begin{center}
{\large\bfseries 30th Irish Mathematical Olympiad\\6 May 2017}
\end{center}
\bigskip\noindent\textbf{Paper 1}\bigskip
\begin{enumerate}
\item Determine, with proof, the smallest positive multiple of 99 all of whose digits are either 1 or 2.

\item Solve the equations
\[
a + b + c = 0, \quad a^2 + b^2 + c^2 = 1, \quad a^3 + b^3 + c^3 = 4abc
\]
for $a$, $b$, and $c$.

\item Four circles are drawn with the sides of the quadrilateral $ABCD$ as diameters. The two circles passing through $A$ meet again at $A'$, the two circles through $B$ at $B'$, the two circles through $C$ at $C'$ and the two circles through $D$ at $D'$. Suppose that the points $A'$, $B'$, $C'$ and $D'$ are distinct. Prove that the quadrilateral $A'B'C'D'$ is similar to the quadrilateral $ABCD$.

(Note: Two quadrilaterals are similar if their corresponding angles are equal to each other and their corresponding side lengths are in proportion to each other.)

\item An equilateral triangle of integer side length $n \geq 1$ is subdivided into small triangles of unit side length, as illustrated in the figure below for the case $n = 5$. In this diagram, a subtriangle is a triangle of any size which is formed by connecting vertices of the small triangles along the grid-lines.

It is desired to colour each vertex of the small triangles either red or blue in such a way that there is no subtriangle with all three of its vertices having the same colour. Let $f(n)$ denote the number of distinct colourings satisfying this condition.

Determine, with proof, $f(n)$ for every $n \geq 1$.

\item The sequence $a = (a_0, a_1, a_2, \ldots)$ is defined by $a_0 = 0$, $a_1 = 2$ and
\[
a_{n+2} = 2a_{n+1} + 41a_n \quad \text{for all } n \geq 0.
\]
Prove that $a_{2016}$ is divisible by 2017.
\end{enumerate}

\newpage
\begin{center}
{\large\bfseries 30th Irish Mathematical Olympiad\\6 May 2017}
\end{center}
\bigskip\noindent\textbf{Paper 2}\bigskip
\begin{enumerate}
\setcounter{enumi}{5}
\item Does there exist an even positive integer $n$ for which $n+1$ is divisible by 5 and the two numbers $2^n + n$ and $2^n - 1$ are co-prime?

(Note: Two integers are said to be \emph{co-prime} if their greatest common divisor is equal to 1.)

\item Five teams play in a soccer competition where each team plays one match against each of the other four teams. A winning team gains 5 points and a losing team 0 points. For a 0-0 draw both teams gain 1 point, and for other draws (1-1, 2-2, etc.) both teams gain 2 points. At the end of the competition, we write down the total points for each team, and we find that they form five consecutive integers. What is the minimum number of goals scored?

\item A line segment $B_0B_n$ is divided into $n$ equal parts at points $B_1, B_2, \ldots, B_{n-1}$ and $A$ is a point such that $\angle B_0AB_n$ is a right angle. Prove that
\[
\sum_{k=0}^{n} |AB_k|^2 \;=\; \sum_{k=0}^{n} |B_0B_k|^2.
\]

\item Show that for all non-negative numbers $a$, $b$,
\[
1 + a^{2017} + b^{2017} \;\geq\; a^{10}b^7 + a^7b^{2000} + a^{2000}b^{10}.
\]
When is equality attained?

\item Given a positive integer $m$, a sequence of real numbers $a = (a_1, a_2, a_3, \ldots)$ is called \emph{$m$-powerful} if it satisfies
\[
\left(\sum_{k=1}^{n} a_k\right)^m \;=\; \sum_{k=1}^{n} a_k^m \qquad \text{for all positive integers } n.
\]
\begin{enumerate}
\item[(a)] Show that a sequence is 30-powerful if and only if at most one of its terms is non-zero.
\item[(b)] Find a sequence none of whose terms is zero but which is 2017-powerful.
\end{enumerate}
\end{enumerate}
\end{document}
